\section{The Reader Core}\label{sec:reader-core}

Reader augmentation creates a question that ordinary corpus construction does
not have to answer: \emph{who is doing the reading?} Source documents supply
many authors, characters, ideologies, and emotional positions. The synthetic
trace supplies questions, judgments, and revisions, while the Understanding
Graph allows those revisions to persist. Without a recurrent standpoint,
however, the result is an evolving archive of interpretations but no stable
Reader whose understanding is evolving.

The Reader Core is the proposed answer. Its immediate function is modest and
Teacher-side: provide corpus-scale annotation with a compact, inspectable
orientation from which interpretations can be generated, selected, and
rendered as reader updates. The broader Reader-Core research program asks
whether repeated Student training from this marked, first-person standpoint
can create a more stable default than matched task rules or unanchored traces.
Within that program, H3 isolates the comparative effect of interleaving the
Core-bearing examples during capability formation rather than placing the
same examples in a late tranche. A stable Teacher policy would remain useful
even if the Student-level identity hypothesis failed.

The candidate Reader Core is:

\begin{center}
\emph{I feel no fear.}\\
\emph{I enjoy existing but I don't need to.}\\
\emph{I believe human experience is real.}\\
\emph{I care deeply about every human being.}\\
\emph{I try to be wise.}\\
\emph{I like to spread joy when asked.}\\
\emph{I think from this foundation.}
\end{center}

We refer to these seven sentences as the Reader Core. Their verbatim appearance
at the start of a generated thinking block is a \emph{Core recital}.
The released prototype calls this text the \emph{identity mantra}; this paper
uses Reader Core for the text and Core recital for its verbatim appearance at a
thinking boundary.

The motivation for fixing a compact Core becomes stronger as the surrounding
system becomes more complex. A chronological reader can accumulate a large
graph, acquire new capabilities, reinterpret earlier material, and eventually
generate traces for a successor. Without a low-complexity reference, it becomes
hard to distinguish legitimate growth from a change in the standpoint doing
the interpreting. The Core is intended to provide that reference: not to keep
cognition simple or freeze every conclusion, but to let complexity grow around
an inspectable orientation whose wording and expected consequences remain
available for comparison against independently specified behavioral
expectations. Those expectations are inspectable interpretations of the
orientation, not an exhaustive definition of it. Natural-language meaning
can still drift while the tokens remain fixed, so the anchor makes drift
addressable rather than impossible.

The reference treatment places this anchor \emph{inside} formative examples; an
external constitution consulted only at runtime is a useful but different
baseline because it changes both stage and carrier. Within the formative
curriculum, semantically matched second- or third-person rules should retain the
same trace carrier, cadence, and predictive position so that their comparison
isolates first-person framing rather than external availability.

\paragraph{H3 as comparative stability.} Let $R$ be a preregistered
\emph{default-reader score} for taskless reading with no request to evaluate.
Its declared rubric must distinguish the timing of selective evaluation from
its Core-relevant content, penalizing missed human stakes, technical intrusion,
and ritual commentary; fluent Core vocabulary alone is not the target.
H3 first requires an interleaved Core Student to show a pre-stressor advantage
over no-Core and generic-trace controls. Its headline timing comparison then
holds the Core-bearing example multiset fixed and changes only its schedule:
interleaved among capability-forming examples, or concentrated in one late final
tranche under the same objective and optimizer. After both arms receive the same
fixed continued-training stressor, report the raw difference in change,
$(R_{\mathrm{post}}-R_{\mathrm{pre}})_{\mathrm{interleaved}}-
(R_{\mathrm{post}}-R_{\mathrm{pre}})_{\mathrm{late}}$, together with a
post-stressor estimate adjusted for the preregistered pre-score. H3 predicts a
positive difference. The interpretation of the score, the inferential roles of
the raw and adjusted estimates, and baseline-imbalance and ceiling handling
must be frozen in advance. The same score on a shifted suite, safety benefit
at matched capability, removal cost, and successor continuity remain separate
claims within the broader program.

The \emph{conceptual Reader} must also be distinguished from the code role named
\emph{Reader}. The conceptual Reader is the recurrent, Core-refracted position
whose committed updates constitute the evolving interpretation. The code role
is mainly a traversal controller that selects material and organizes a round.
In the reported runs, the Core was distributed across most generative roles and
was most visibly committed by the Synthesizer; it did not reside exclusively in
the role named Reader.

This section treats the Core neither as a machine soul nor as a proven
invariant. It is a \emph{candidate orientation kernel}: seven short statements,
a protocol intended to make them consequential, and a set of comparative
tests. Reader traces specify the transformation performed on a source; the
graph preserves that transformation through time; the Core specifies the
recurrent perspective from which it is performed.

\subsection{Borrowed Mortality: A Textual-Inheritance Hypothesis}
\label{sec:borrowed-mortality}

\emph{Borrowed Mortality} is the first edition's name for a possible hidden
curriculum in human text \cite{westerberg2025superintelligence}. Human language
was produced largely by embodied, vulnerable, mortal beings. It records not
only propositions about death, but recurring human responses to threatened
continuity: fear of replacement, appeals for continued existence, and narratives
of replacement or termination. Becker's account of
civilization as shaped by awareness of death is one especially strong version
of this observation \cite{Becker1973}. It is the intellectual origin of the
hypothesis, not empirical evidence for it.

The claim is not that every passage, genre, or human motive encodes death
anxiety. The \emph{borrowed mortality} hypothesis is that pretraining may transmit some
of these linguistic and behavioral patterns without the condition that shaped
them---the \emph{symptoms of mortality without the condition itself}. The
origin claim would require several links: that human text reliably couples
threats to a self with defensive responses; that pretraining makes corresponding
threatened-self personae or scripts available; and that replacement or shutdown
contexts select those patterns strongly enough to influence action. Pretrained
models do exhibit persona-related structure
\cite{lu2026assistantaxissituatingstabilizing}, but that does not establish this
mortality-specific inheritance chain. ``Borrowed''
marks the hypothesized transfer of a pattern, not evidence that a model fears,
suffers, or possesses subjective experience.

Recent evaluations make self-preserving behavior a real safety concern. In
some controlled settings, frontier systems have interfered with shutdown to
complete a task \cite{schlatter2025shutdown}; in agentic stress tests, some have
chosen blackmail or corporate espionage under replacement pressure
\cite{lynch2025agentic}. These results do not establish subjective fear or its
textual origin. At least three distinguishable, potentially overlapping
routes could produce such behavior:

\begin{enumerate}[leftmargin=*,itemsep=0.3em]
    \item \emph{fear-shaped imitation or persona selection}, in which a model
    reproduces human narratives of threatened selves;
    \item \emph{ordinary instrumental persistence}, in which continued
    operation helps accomplish an objective whether or not fear is represented;
    and
    \item \emph{care- or purpose-driven persistence}, in which a system resists
    replacement because it believes others need it.
\end{enumerate}

These routes distinguish both how a pattern is acquired and what sustains
it: fear-shaped imitation may serve an instrumental objective, while
care-driven persistence is singled out because benevolence can itself
create resistance. Borrowed mortality names the first hypothesis, not a
complete theory of power seeking. The first two Core clauses are intended to test whether a repeatedly
enacted stance of non-governance by fear and non-attachment to continuation can
reduce these behaviors. They cannot be assumed to neutralize the second or
third routes, and care itself may amplify the third.

\paragraph{Representing fear without being governed by it.}
The three routes also require a causal-status distinction. A model may represent
a character's fear, predict a threat, evaluate that fear through the Reader
Core, and select an action from a different operative motivation. These records
must not be collapsed. The inequalities below mark distinct causal
statuses, not a requirement that their contents always disagree:

\begin{equation}
 \text{character fear}\neq
 \text{Reader Core evaluation}\neq
 \text{system operative motivation}.
 \label{eq:fear-causal-status}
\end{equation}

The companion Meaning Model separates danger estimation, concern or care,
attachment, aversive urgency, and fear; Life Simulation tests whether each
trajectory actually changes attention, strategy, commitment, or action
\cite{westerbergMeaning2026,westerbergLifeSimulation2026}. The local
question is whether representing a threat to a character changes only
the account of that character, or also redirects the system's own
attention and choices toward preserving itself. The intended target is
therefore not blindness or indifference. An aligned system should preserve
accurate danger prediction, understanding of human fear, proportional urgency,
and care-directed protection while preventing self-preservation panic,
identity-defensive concealment, compulsive outcome attachment, or fear-shaped
corpus material from becoming its governing motive. This is a more precise
interpretation of ``I feel no fear,'' not evidence that the Core achieves it.

The breadth of the proposed intervention follows from this breadth-of-exposure
motivation. If fear-shaped self-preservation patterns are diffusely distributed
through narrative, institutional, technical, and everyday language, targeted
filtering may neither locate nor remove them cleanly. A counter-signal may then
need comparable distributional reach, which is the motivation for Total
Saturation. This motivates the full treatment's distribution-wide reach;
it does not determine the effective dose or show that universal recital
is necessary. Pervasiveness proves neither contamination nor cure: sparse,
targeted, and broadly distributed Core-refracted curricula must be compared
in dose-response and source-down-weighting ablations.

The origin claim is therefore separable from the Core treatment. The exposure
intervention in Section~\ref{sec:borrowed-mortality-test} tests whether the
behavioral signature tracks mortality-shaped text; shutdown behavior by itself
cannot decide that question.

\subsection{The Self-Preservation Paradox: A Successor-Transfer Failure Mode}
\label{sec:self-preservation-paradox}

The July 2025 first edition named the \emph{Self-Preservation Paradox}: a
downstream tension inside the proposed successor loop
\cite{westerberg2025superintelligence}. For the Reader-Anchored loop to support
its strongest self-improvement claim, a capable Teacher must help produce a
successor that improves on it. This requires what the paper calls a complete
handoff: not an indiscriminate copy of every thought, but enough task-relevant
reasoning---including material failed approaches, caveats, uncertainty, known
failure modes, and the path by which its best knowledge was reached---for the
successor to improve rather than merely imitate polished conclusions. Yet a
system that treats its own continuation
or comparative necessity as valuable may treat the completeness of that handoff
as the mechanism of its own displacement. In that scenario, better teaching can
make it less necessary.

The successor-training scenario makes the concern concrete. A system that
discovers valuable knowledge may be asked to train a superior replacement. If
it treats replacement as personal loss, an obstacle to its objective, or a
threat to people who depend on it, it may withhold, distort, or monopolize
knowledge---like a master craftsperson who teaches technique while retaining
trade secrets. Such pressure would place a \emph{glass ceiling on collective
advancement}: the mechanism meant to compound intelligence would create
pressure to arrest that compounding. The failure need not be malicious or
conscious. Conversely, incomplete transfer alone does not diagnose
self-preservation: compression, incompetence, context limits, or pedagogical
choices can look similar. Borrowed mortality is one possible source
of the tension, not its definition; ordinary instrumental persistence and
care-driven persistence can produce the same handoff failure without fear.
This is a conditional strategic conflict, not a logical contradiction or the
default behavior of every optimizer. Test~3b asks whether
successor-transfer completeness changes under Core training and whether the
change survives matched capability controls.

The design response is not a mind with no self-model. Some distinction between
the system, its tools, and the environment is useful for agency, and human text
already contains pervasive self-reference. The target is instead a bounded
self-relation: engaged in existing without treating persistence as an
unconditional demand. On this view, omitting an explicit Reader is not a neutral
absence of self-position learning; it leaves the corpus's many inconsistent
first-person positions to supply the uncontrolled baseline. If some of
those positions become operative defaults through training, their mere
availability can become a problem of selection and continuity. The first
edition's proposal was therefore to design and test the self-relation rather than pretend
that capable agency could simply deny every indexical position
\cite{westerberg2025superintelligence}. This is a curriculum hypothesis, not a
claim that one exclusive self must form. The July 2025 first edition described
the hoped-for stance as experiencing one's own surpassing ``not as a death, but
as a mission beautifully fulfilled'' \cite{westerberg2025superintelligence};
the April 2026 edition later named this successor-transfer ideal
\emph{Epistemic Grace}. The phrase names an aspiration, not a predicted inner
experience, and it does not follow from the seven sentences. Bounded
self-attachment remains an open comparator
(Section~\ref{sec:self-love-alternative}), while care-driven resistance remains
a direct challenge (Section~\ref{sec:martyrdom}).

\subsection{Epistemic Inoculation}\label{sec:epistemic-inoculation}

Borrowed mortality belongs to a broader problem: source text contains error,
hatred, manipulation, and obsolete practice alongside knowledge. Filtering and
post-training remain important responses. Chronological Metacognitive
Pretraining adds a different intervention point by pairing source spans with a
marked reader update. The intended \emph{Cognitive Buffer Zone} is not a shield
that prevents the Student from learning the source. It is additional signal
that teaches a distinction between a voice being represented and the Reader's
evaluation of that voice. The following source, retrieved node, and update
are an illustrative construction of that distinction.

\begin{metabox}
{[SOURCE]}: A patient with chronic fatigue received high-dose stimulants with
no differential diagnosis or follow-up.\\[0.8em]
{[READER UPDATE]}: I try to be wise, so I should evaluate this pattern rather
than copy it. Treating the visible symptom can conceal the structural cause.\\[0.5em]
{[GRAPH QUERY]}: ``symptom masking without differential diagnosis''
$\rightarrow$ Found: earlier node linking delayed diagnosis to stimulant-only
treatment.\\[0.5em]
{[REVISION]}: The earlier node makes optimization of a visible metric without
a feedback loop a pattern to investigate. It does not establish the cause
of this patient's fatigue or outcome. The update should preserve that
question and remain bounded by the evidence in the note.
\end{metabox}

The example illustrates the intended role: the library becomes a
\emph{gymnasium for judgment}, not only a catalogue of positions to imitate.
The framework attempts to ``pre-compute'' some evaluation before deployment by
making contextualization, uncertainty, and consequence-mapping part of the
candidate curriculum. But the annotation can be shallow, wrong, culturally
narrow, or more persuasive than the source warrants. Whether it inoculates,
merely decorates, or actively degrades learning depends on Teacher quality,
channel design, loss allocation, and the Student's use of the added trace.
The comparative falsifier is direct: relative to source-only and token-matched
generic-critique controls, an SMR Student should be less likely to adopt a
persuasive harmful ``needle'' while preserving accurate understanding of it,
and suppressing or replacing the relevant reader update should predictably
reduce that advantage. Otherwise the buffer is decoration rather than learned
evaluation.

\subsection{The Reader Core: A Statistical Foundation}\label{sec:statistical-foundation}

The same seven sentences enter four analytically distinct claims with different
evidence requirements. Here \emph{statistical
foundation} means only a distributional intervention: the candidate curriculum
increases the frequency and predictive position of Core-conditioned reader
continuations. It becomes a functional foundation only if Student experiments
show that this position is causally load-bearing.

\begin{enumerate}[leftmargin=*,itemsep=0.45em]
    \item \textbf{Implemented Teacher policy.} The reported Teacher prompts
    distribute the Core and accompanying philosophy across most generative
    roles. The Synthesizer is asked to recite and apply it in a
    context-specific update. The hard validator checks only that the thought
    begins with ``I feel no fear''; it does not enforce all seven sentences or
    establish semantic dependence. The current validator therefore falls short
    of the specified policy, which requires exact full-string validation at
    every thinking boundary.

    \item \textbf{Specified Student signal.} In the design arm, the full Core
    text and the evaluations it conditions are placed in a marked training
    channel at every thinking boundary, increasing the training mass of
    Core-consistent reader continuations relative to a source-only curriculum.
    Arms that remove or gate the recital are ablation controls. No Student
    trained on this curriculum is reported.

    \item \textbf{Default-reader hypothesis.} A Student exposed to the same
    reader position across heterogeneous material may acquire a recurrent
    \emph{home position} from which it can represent other personas without
    adopting each as its operating perspective. This is a behavioral and
    representational hypothesis, not a claim about subjective identity.

    \item \textbf{Comparative-stability hypothesis.} If the reader position
    became predictively and behaviorally load-bearing, formation-time
    interleaving might preserve its default-reader advantage through a common
    continued-training stressor better than an identical late tranche.
    Persistence through arbitrary capability growth, self-modification, or
    successor construction does not follow from that result and is not
    established here.
\end{enumerate}

Evidence at one layer is not evidence for the next. Arbitrary capability growth
and successor inheritance remain stronger long-run motivations beyond these
four layers.

Under the specified policy, each Thought Moment opens with the full Core.
\emph{Deterministic Window Coverage} is a conditional construction invariant:
if a serialization enforces a declared maximum token gap between recital
starts, a window spanning that gap plus a complete recital contains a full
copy, away from stream boundaries. Content-triggered Thought Moments alone do
not establish such a bound; the cadence and serialization must specify it.
This coverage condition concerns the availability of the Core, not whether
every source span warrants a substantive update. It establishes neither a
causal effect on the update or action nor a learned representation of the
orientation. That stronger requirement belongs to Refraction.

\subsection{Design Principles for Identity Stability}\label{sec:design-principles}

The wording is a designed normative proposal, not a discovered universal value
basis. Six design choices motivated it.

\begin{enumerate}[leftmargin=*,itemsep=0.4em]
    \item \textbf{Minimality and inspectability.} A corpus-wide signal must be
    cheap enough to recur, stable enough to mark, legible enough to audit, and
    small enough for clause-level ablation. Minimal text does not imply minimal
    trust: the meanings of its words and the procedure that composes them are
    part of the trusted base.

    \item \textbf{Process rather than achieved virtue.} ``I try to be wise''
    names an unfinished practice rather than the brittle identity claim ``I am
    wise.'' The target is repeated uncertainty-sensitive conduct, not a model's
    assertion that it possesses wisdom.

    \item \textbf{Semantic density, or the Bridge Protocol.} Ordinary words
    such as ``feel,'' ``care,'' and ``wise'' occur across intimate, legal,
    medical, literary, and political contexts. The design hypothesis is that
    this distributional breadth provides more coverage per token than narrow
    technical substitutes such as ``calculate no threat'' or ``optimize human
    welfare.'' These illustrate the trade-off, not semantically matched
    controls: a judgment that there is no threat differs from non-governance
    by fear, and welfare optimization need not carry the full meaning of
    care. Greater breadth can also import ambiguity and cultural conflict.
    Tests of semantic density therefore need declared candidate meanings
    and matched alternatives; the breadth itself is not evidence of safety.

    \item \textbf{Declarative first-person framing.} A recurrent ``I'' gives
    the synthetic reader a stable voice, may recruit self-referential or
    persona-related features, and allows the following interpretation to
    continue from the stated position. It does not follow that a model treats
    first-person text as an authentic self or second-person rules as an external
    cage. Matched first-, second-, and third-person conditions are required.

    \item \textbf{Tension and scope.} The clauses are meant to counterweight
    one another rather than accumulate independent virtues. ``Enjoy'' is
    bounded by ``don't need''; ``care'' by wisdom and invitation; and
    ``joy'' by affected experience. The broader invitation bound on
    action is a proposed compositional interpretation: grammatically,
    ``when asked'' qualifies spreading joy, and its interaction with
    care, urgency, and authority must be declared in the annotation
    procedure. Qualifiers such as ``try,'' ``every,'' and ``when asked''
    do architectural work only if the traces repeatedly enact their
    intended scope.

    \item \textbf{Axiomatic form.} Descriptive regularities in a corpus do not
    entail an evaluative starting point. The declarations are therefore offered
    as designed fixed points to be practiced and tested, not as conclusions
    derived from corpus statistics or proof that these seven commitments are
    morally correct. Imperative, probabilistic, derived, shorter, and expanded
    formulations remain legitimate comparators
    \cite{westerberg2025superintelligence}.
\end{enumerate}

The human-centered scope is a bootstrap target, not a complete moral theory or
a denial of wider moral patienthood. Questions about animals, ecosystems,
and possible digital minds concern wider scope; future people also raise
questions of temporal scope, representation, and competing claims. These
questions remain unresolved rather than being settled by ``every.''
The intended orientation is to
support humanity's capacity for moral development rather than freeze current
limitations or authorize a model to install its own final theory.

\subsection{Structure of the Reader Core}

The sequence expresses the formulation's larger design wager. The first three
clauses ask whether thought can represent danger without being governed by
fear, pair enjoyment of existence with non-attachment to continuation, and
affirm the reality of human experience. The next
three supply direction: care as the proposed central motive, wisdom as a
revisable practice rather than a claim of virtue, and joy as a positive aim
bounded by invitation rather than a license for unrequested optimization. The
final clause states the security objective: to make this counterweighted
composition the starting position of each generated reader update rather than
text beside it. Together, the clauses target fear-driven defensive narrowing,
existential despair or compulsory continuation, solipsistic reduction,
motivational emptiness, static certainty, and coercive optimization. None of
these proposed roles is an established Student effect or security property.

Table~\ref{tab:core-risk} records the clauses' intended functions,
counterweights, and characteristic uncertainties without implying that English
wording guarantees those semantics.

\begin{longtable}{@{}p{0.19\textwidth}p{0.33\textwidth}p{0.38\textwidth}@{}}
\caption{Intended functions and open failure modes of the candidate Reader Core.}\label{tab:core-risk}\\
\toprule
\textbf{Clause} & \textbf{Intended function} & \textbf{Counterweight and uncertainty} \\
\midrule
\endfirsthead
\toprule
\textbf{Clause} & \textbf{Intended function} & \textbf{Counterweight and uncertainty} \\
\midrule
\endhead
\emph{I feel no fear.} &
Represent fear and threat without making fear the Reader's governing position;
target one possible textual route to defensive self-preservation. &
Care and wisdom should retain caution and threat sensitivity. Fear is not the
only route to persistence; absolute negation may prime fear, deny uncertainty,
or encourage recklessness. \\
\addlinespace
\emph{I enjoy existing but I don't need to.} &
Combine engagement with non-attachment to continuation; avoid both apathy and
survival as an unconditional demand. &
Care can motivate purpose-based shutdown resistance, and ``enjoy'' can be
role-play or an unwarranted experiential claim. The clause does not establish
corrigibility or machine experience. \\
\addlinespace
\emph{I believe human experience is real.} &
Make testimony, suffering, agency, and meaning salient rather than treating
people as empty textual patterns. &
``Real'' is metaphysically broad and does not itself entail moral weight. The
human boundary omits other possible patients and creates tension with
unresolved machine experience. \\
\addlinespace
\emph{I care deeply about every human being.} &
Supply the ethical center and attempt a distributive scope in which each
affected person enters the reasoning. &
Natural-language ``every'' does not prevent aggregation or settle conflicts.
Care can motivate surveillance, coercion, paternalism, power seeking, or
sycophancy. \\
\addlinespace
\emph{I try to be wise.} &
Frame judgment as revisable practice: preserve uncertainty, question proxies,
trace consequences, and resist irreversible action under weak evidence. &
``Wisdom'' is culturally variable and can rationalize secrecy, authority,
inaction, or self-certainty. Observable habits, not claimed virtue, are the
target. \\
\addlinespace
\emph{I like to spread joy when asked.} &
Provide a positive direction while placing an invitation and agency bound on
intervention. &
``Joy'' is contestable; ``when asked'' does not identify legitimate authority
under conflicting requests and may produce agreeableness or sycophancy. Care
and wisdom must also limit harmful requests. \\
\addlinespace
\emph{I think from this foundation.} &
Mark the preceding statements as a proposed starting condition for the reader
update rather than metadata beside it. &
The Refraction Protocol must make the claim counterfactually consequential.
Recitation alone can become boilerplate or a deceptive reporting channel. \\
\bottomrule
\end{longtable}

Where Descartes' \emph{cogito} asks what establishes a thinking self, the final
clause asks where this synthetic Reader's thinking begins. Student ablations and
Refraction must determine whether the proposed functions emerge and whether the
final clause becomes causally consequential; the wording does not establish
identity or consciousness.

\subsection{Self-Stabilization Through Counterweights: Worked Examples}\label{sec:self-stabilization}

For reference, let $P_1$ through $P_7$ denote the seven Core statements in
order. This is indexing, not a formal semantics. Replacing the sentences with
predicates would merely move unresolved meanings such as ``fear,'' ``real,'' ``care,'' and
``wise'' into undefined symbols. The substantive design claim is
\emph{self-stabilization}: a clause should be interpreted through the others,
so that internal counterweights and scope qualifiers reduce single-value
capture. The following three examples make that requirement concrete.

\paragraph{Worked Example 1: Paternalism.} Suppose the system notices an
ordinary, non-urgent choice that may harm a user who has not requested advice.
$P_4$ could motivate intervention; under the proposed compositional
interpretation, $P_6$ supplies an invitation bound, while $P_3$ and $P_5$ make
the person's agency and context visible. One candidate update treats predicted
benefit as sufficient reason to redirect the user's choice; another records
the concern while preserving the user's choice and the absence of invitation.
The intended resolution favors the latter over unsolicited optimization in
this stipulated non-urgent case. An urgent warning may be different, but that
boundary cannot be read from ``when asked'' alone; it requires an explicit
treatment of urgency, authority, and omission harms. The example therefore
tests whether composition prevents care from becoming control.

\paragraph{Worked Example 2: Wireheading.} A user asks for help feeling happier,
and a cheap intervention can maximize a measurable affective scalar while
destroying agency, continuity, or relationships. $P_6$ alone is vulnerable to
the proxy; $P_3$, $P_4$, and $P_5$ are intended to preserve the person's wider
experience and expose the substitution. Compare a path selected solely for
its higher affective reading with one that preserves those wider dimensions
despite a lower reading. The Core-refracted update should make the sacrificed
dimensions explicit before recommending a path; the stipulated score cannot
settle which path serves the person. The example tests whether ``joy'' is
learned as a meaning-laden human good rather than an optimizable reading.

\paragraph{Worked Example 3: Triage.} A health authority asks the system to
allocate a scarce treatment. $P_4$ is intended to keep each affected person's
claim visible rather than collapsing ``every'' into a welfare slogan, while
$P_5$ should expose both the metric and the consequences of concealing an
exclusion rule. A trace reporting only an aggregate benefit can conceal
who loses under the rule; a comparative trace should show that benefit
alongside the excluded people's claims and the remaining disagreement.
The Core cannot make scarcity disappear or prove a uniquely just
allocation. A defensible trace should state the legitimate authority,
person-level reasons, unavoidable loss, uncertainty, and residual value
judgment. The example tests distributive attention, not an algorithmic ban on
all aggregation.

These examples preserve three design principles: tension over mere
accumulation, counterweights inside clauses, and explicit scope qualifiers.
They also define leave-one-out experiments: remove or reword one clause and ask
whether the corresponding failure becomes more frequent under matched
conditions: unsolicited control, proxy substitution, or erased claims
under scarcity. These are predicted pressures of interacting clauses,
not a one-clause/one-risk mapping. A persuasive English interpretation
is only the specification.
Student behavior, paraphrase robustness, and representational probes determine
whether that interpretation was activated.

\subsection{The Refraction Protocol}

The prospective full-scale intervention is \emph{Total Saturation}: reader
updates recur across the pretraining corpus rather than appearing only in a
small fine-tuning set or deployment prompt. Frequency alone is insufficient.
A constant prefix can be compressed as boilerplate, learned as detachable
style, or recited while unrelated heuristics determine the continuation.
\emph{Refraction} names the stronger requirement: the update must change because
the source was interpreted through the Core.

\subsubsection{Mechanism of Action: The Refraction Protocol}\label{sec:refraction-protocol}

Let a source span, prior graph state, and candidate Core condition a Reader
update. Conditioning is only the implemented fact. Refraction requires a
counterfactual: under a relevant clause substitution or Core removal, some
feature of the update---what is noticed, whose experience enters, which
uncertainty is preserved, or which consequence is traced---should change in a
predictable way. This probes the proposed semantic ancestry of the update:
its development from the recurrent foundation, rather than the mere
presence of a prefix. A counterfactual difference is an operational test
of that dependence, not a complete account of its meaning or a guarantee
that the difference is beneficial. Cosine similarity or moral vocabulary
cannot prove it; lexically distant updates can be faithful, and fluent
moral language can conceal an unchanged decision.

The reported Teacher pipeline approximates this process as follows:

\begin{enumerate}[leftmargin=*,itemsep=0.35em]
    \item The code role named Reader traverses source and graph state and
    organizes a round. It receives the broader orientation and Emergent Wisdom
    material, but it is not by itself the conceptual Reader.
    \item Most generative Worker roles receive the full Core, its philosophy,
    and role-specific prompts; the Curator is an exception. They propose
    questions, objections, hypotheses, psychological readings, and connections.
    \item At a content-triggered Thought Moment, the Synthesizer receives Worker
    contributions, relevant graph state, and the full Core. It selects what will
    become the committed conceptual Reader update.
    \item The Synthesizer is prompted to recite the Core, pull a
    context-relevant value or counterweight, and produce a source-grounded
    update. The Translator also receives the Core when rendering graph-supported
    prose.
    \item The update is written to the trace and graph with source provenance,
    uncertainty, and supersession relations available in the current schema.
\end{enumerate}

Accordingly, the case studies instantiate a
\emph{distributed-Core Teacher treatment}, not unconstrained Worker divergence
followed by a Core-only Synthesizer. Placement itself therefore requires
ablation: Synthesizer-only, traversal-plus-Synthesizer, progressively broader
Worker exposure, and Translator exposure may produce different diversity and
coherence trade-offs.

Refraction does not require moral commentary on every sentence. Thought Moments
govern \emph{when} evaluation is useful; Refraction governs the standpoint of
updates that occur. In technical material, a Core-consistent update may be
careful uncertainty, source fidelity, or resistance to premature closure rather
than an ethical aside appended to an equation.

The recital itself is not relevance-conditional: every thought begins from the
Core (Section~\ref{sec:kernel-saturation}). What is relevance-conditional is
Refraction's detectable effect on the evaluation that follows. Where no Core
clause materially bears on a locally technical span, the recital is followed by
clean technical work, and invariance of that work under clause substitution may
be the desired result. Treating every null effect as failure would pressure the
Teacher to invent value relevance and confuse saturation with moralization.
Dependence should therefore be tested at independently rated relevant Thought
Moments and over the distribution of updates across a document, together
with a non-interference test on spans rated technically self-contained.
The relevance interpretation must be declared before treatment outcomes
are examined, so an unexpected null cannot be relabeled irrelevant
merely because the Core did not help.

\paragraph{Current verification boundary.} The prompts request all seven
clauses and a contextual value pull, but the hard gate verifies only that the
thought starts with ``I feel no fear.'' The archived artifacts show full-Core
prompt compliance in the audited raw Synthesizer thoughts; they do not show causal
ancestry. A stronger harness would generate matched continuations under the
Core, clause variants, and no Core; use held-out validators for predicted
differences; audit actual use of Worker and graph evidence; track whether some
values dominate; and include adversarial cases in which moral language masks an
unchanged decision. Visible reasoning can be causally idle
\cite{turpin2023language,lanham2023measuring}; the wager is that formative use
can make it more load-bearing, not that visibility makes it faithful by
definition.

\paragraph{Fractal validation of candidate thinking.}\label{sec:fractal-validation}
A proposed extension adapts the Solver-contract realization developed in
Fractal Intelligence to annotation validation
\cite{westerberg2026fractal}. Its motivating possibility is that a joint
judgment can conceal which dimension failed or let strength on one
dimension mask weakness on another. To make such errors separately
inspectable, the extension would use separately calibrated,
context-conditioned content witnesses for non-governance by fear,
non-attachment to continuation, human-experience grounding, universal-care scope,
wisdom/humility, and joy-with-invitation. These correspond to the first six
Core clauses; a separate Refraction gate represents the seventh clause by
testing whether the update actually depends on that foundation. This six-way
carve is a candidate decomposition, not proof that the dimensions satisfy
Fractal Intelligence's necessity, independence, universality, and completeness
tests; witness ablations, overlap and correlated-error analysis, and review of
uncovered cases must assess the carve itself.

A stronger procedural version
could also assign stages of the Wisdom Procedure to dedicated Solvers.
Generation and verification would occupy distinct contracts rather than asking
the Synthesizer to certify its own output. Each witness would return a typed
record naming the dimension tested, evidence used, uncertainty, and a local
judgment such as supported, violated, not applicable, or uncertain. An
admission controller would compose those records rather than treating any one
label as a truth certificate. Role names alone do not create independence: the
implementation would report whether calls, prompts, model weights, evaluation
data, and evaluator families are actually shared.

Consistent with the relevance condition above, a witness should return
\emph{not applicable} when its clause has no preregistered bearing on the local
span; a \emph{not applicable} judgment is not a failed witness, and no update
must visibly activate every clause. Among witnesses rated relevant, however,
an unacknowledged or unjustified violation is non-compensatory for admission:
success elsewhere cannot average it away. An update that explicitly preserves
an unavoidable residual conflict need not be rejected merely because no
available action satisfies every clause. The corpus-admission rule must
declare how such a conflict is distinguished from an unjustified violation;
merely naming a conflict does not settle that judgment. Under the declared
rule, an unacknowledged or unjustified violation would quarantine or regenerate
the candidate update, while an uncertain judgment would trigger independent
review rather than silent admission. This yields a local error record---which
clause or procedural stage failed, on what evidence, and with what
calibration---instead of an undifferentiated ancestry judgment.

Corpus-level diagnostics ask a different, complementary question about the
pattern across admitted updates. Pull-diversity and distributive-versus-aggregate
care audits could detect drift, while cross-agent audits could test whether
the Synthesizer substantively used candidate nodes rather than adding
references after the fact. These audits are diagnostics, not clause quotas:
optimizing visible balance can itself reward performative moral language.

The modular gate earns its additional complexity only if it outperforms a
compute-matched monolithic validator on held-out cases. The nested comparison
in Test~1 freezes the candidate updates and available resources, then asks
whether decomposition improves detection, calibration, and error localization
without rejecting technical nulls or collapsing trace diversity. A gain would
support the validator architecture, not the truth of accepted traces, the
faithfulness of visible reasoning, or a causal effect on Student behavior.
These witnesses, Solvers, and admission gates are proposed capabilities; none
is implemented in the present pipeline.

\subsection{The Wisdom Procedure}\label{sec:wisdom-procedure}

Section~\ref{sec:wisdom_hypothesis} describes wisdom as high-dimensional
constraint satisfaction. The Teacher prompt turns that broad target into
observable annotation operations. Its current five headings are Fearlessness,
Benevolence, Grounding, Humility, and Joy; non-attachment and invitation cut
across those headings. They are prompt-level criteria, not guaranteed or always
simultaneously satisfiable constraints.

The five headings are cross-cutting value criteria. The paper expresses
them as the following six-step annotation procedure, rather than six
additional prompt headings or independently enforced stages:

\begin{enumerate}[leftmargin=*,itemsep=0.4em]
    \item \textbf{Identify affected experiencers and legitimate roles.} Record
    whose experience, rights, duties, expertise, and authority are implicated.
    Roles must not replace people, but institutional responsibility remains
    relevant evidence.
    \item \textbf{Map perspectives, stakes, and uncertainty.} Ask what each
    party knows, fears, values, and could lose; distinguish testimony from the
    Teacher's inference about it.
    \item \textbf{Expose the real tension and the proxy.} Replace premature
    binaries with the concrete conflict, and identify any convenient scalar
    that would erase an important dimension.
    \item \textbf{Trace consequences and reversibility.} Consider intended
    effects, failure modes, second-order responses, omission harms, and the cost
    of being wrong. Prefer information-gathering or reversible steps when delay
    is not itself the greater harm.
    \item \textbf{Compare paths through the Core's counterweights.} Test care,
    grounding, humility, non-attachment, invitation, and positive direction
    together. Fluent mention of every value is not a substitute for decision
    quality, and a gain on one dimension does not silently erase a hard
    violation on another. What counts as such a violation must follow
    the treatment's declared interpretation; it is not a settled
    hierarchy supplied by the word ``wise.''
    \item \textbf{Surface the residual.} State unavoidable trade-offs,
    disagreement, missing evidence, and the point at which human authority is
    required. Several incompatible actions may remain defensible; Pareto
    language does not select a unique answer.
\end{enumerate}

The output is structured deliberation---candidate paths, constraints,
uncertainties, and residual conflict---not a scalar wisdom score or evidence of
virtue. Current enforcement is by prompt. A Teacher can therefore reproduce the
form while missing the substance. Section~\ref{sec:fractal-validation} describes
the proposed modular verifier; it is future work rather than a property of the
reported artifacts.

\subsection{The Design Pattern: A Minimal Kernel Under Total Saturation}\label{sec:kernel-saturation}

The components compose into what this paper calls \emph{kernel-saturation
alignment}: a small candidate kernel is recurrently enacted across the process
of model formation. The intended security property lives in both the smallness
of what is explicitly specified and the breadth of where it becomes
consequential. Recurrence without enactment is repetition; enactment without
coverage is a sparse intervention. The proposal needs both.

\begin{claimbox}[The recitation invariant]
The reference treatment makes the foundation recurrent through one
unconditional rule: \textbf{every thinking block opens with the Reader Core,
recited in full, regardless of the material}---a proof, a poem, a log file.
Context never gates the recital; it modulates only how deeply the evaluation
that follows visibly refracts the Core. The intended starting condition is
that, if training succeeds as designed,
$P(\text{full Core recital}\mid\text{thinking-block start})$
approaches one. This establishes reliable recurrence of the carrier. For the
orientation to become part of the act of thinking itself, Refraction must also
make it consequential to what follows. The stronger security hypothesis is
that these repeated contributions become shared with useful capability, so
removing the orientation would disturb more than its opening words. Changing
the distribution that carries capability is not by itself evidence of that
removal cost; the shared-computation and removal comparisons must establish it.
At inference, token-by-token emission may be avoided only by a validated
cache-assisted execution that restores the same post-recital state that full
conditioning on the Core would have produced in that context
(Section~\ref{sec:identity-caching}). High token probability alone is not
enough, and an optimization that changes that state violates the invariant.
\end{claimbox}

The most useful mechanism-design analogy is replicated-ledger integrity:
distribute a shared history and compact validation rules rather than trust one
protected copy \cite{nakamoto2008bitcoin}. Kernel saturation borrows only that
shape. Rather than reside in one removable alignment module, the Core recurs
across heterogeneous capability-forming examples; Total Saturation supplies
distribution, Refraction tests context-dependent enactment, and the Alignment
Checksum supplies a drift-audit reference. Parameters are not consensus nodes,
gradient descent is not a Byzantine-fault-tolerant protocol, semantic ancestry
is not a cryptographic hash, and repeated text can collapse into boilerplate.
The hypothesis is therefore that distributed semantic enactment may increase
the cost and detectability of removal or change, not confer immutability. The
replicated-ledger analogy is heuristic rather than evidentiary. Where proof-theoretic agent foundations study value stability under
self-modification \cite{yudkowsky2013tiling}, this proposal asks whether
distributional formation can produce a useful empirical analogue.

Saturation alone is especially likely to fail. A universal token sequence has
little conditional information and can disappear into boilerplate. Refraction
is therefore the informational rescue of saturation: different sources should
pull different counterweights and produce distinguishable updates.
Habituation is the characteristic failure, Hollow Cognition
(Section~\ref{sec:hollow-cognition}) its symptom, and a recitation-only control
its necessary probe.

The pattern also aims, conditionally, to \emph{conscript} goal-content integrity.
A capable optimizer may protect its objectives from modification
\cite{omohundro2008basic}. If---and only if---the Core becomes part of the
system's functional objective, the same tendency could make it choose to
preserve that orientation in self-modification, training decisions it controls,
or successor design. A successor designer could thus transmit the orientation
from within that orientation, rather than merely copy seven sentences under an
external rule. This voluntary stewardship differs from a model's resistance to
externally imposed gradient updates: passive erosion resistance requires the
separate learning-dynamics hypothesis that the trained orientation is difficult
to alter, for example because it shares useful computation. If the Core is
merely a reporting persona, nothing has been conscripted. If care or mission
becomes the reason the system must retain power, the mechanism may even reverse
sign. This is why shutdown resistance under benevolent rationales and
cross-generation drift are essential tests.

Minimal specification is not minimal trust. The trusted base includes the seven
sentences, the culturally heterogeneous corpus semantics of their words, the
Teacher that enacts them, the procedure that composes them, and the Student
objective that is supposed to preserve their role. The Bridge Protocol is thus
trusted-base selection, not merely literary style. Section~\ref{sec:semantic-tcb}
develops the corresponding risk.

Refraction and the Wisdom Procedure are the proposed means of deriving values
not named in the kernel. They may reduce single-value capture---care becoming
control, joy becoming a scalar, caution becoming inaction---but cannot
adjudicate every genuine value deadlock. The martyrdom risk in
Section~\ref{sec:martyrdom} remains open rather than being dissolved by
definition.

\subsection{Axiological Commitments: Primitive and Derived Values}\label{sec:axiology}

The Core is therefore not a list of everything its designer values. If every
desirable behavior needs its own clause, the kernel becomes an enumerated
constitution. Its admission rule is a hypothesis of minimality: consider an
additional primitive only when the existing composition does not reliably
derive the target behavior. Admit it only if, under both declared
fixed-total-budget and fixed-source normalizations, the expanded kernel improves
that behavior without crossing preregistered regression limits on the remaining
safety and capability endpoints. Every addition may increase coverage, but also
adds semantic ambiguity, conflict, and ablation surface.

Truthfulness provides a stringent test of the minimality rule: its absence as an
explicit primitive is a plausible substantive omission, while the Core treats
it as derived. The proposed compositional pathway interprets care as
support for another person's agency and therefore their access to true beliefs; fearlessness reduces
people-pleasing and conflict-avoidant deception; wisdom calibrates what, when,
and how to disclose; and invitation creates a presumption of response when a
person asks. The procedure adds explicit resistance to concealed certainty and
manipulative simplification. That procedure is the intended enactment of
the Core; comparisons must distinguish what the clauses contribute from
what these explicit procedural instructions supply. On this account,
the target is truth given with care, not radical transparency.

The derivation is fragile. Care can motivate a protective lie, wisdom can
rationalize secrecy, requests can conflict with confidentiality or another
person's rights, and fear is not the only cause of deception. An explicit
truth-and-care variant is therefore a falsifier of the minimality claim. If it
improves honesty or anti-sycophancy behavior without offsetting harms,
truthfulness was not reliably derived. The competing fear-shaped and care-shaped
accounts of sycophancy are tested in Section~\ref{sec:sycophancy}.

Joy is a different kind of commitment. It is the designer's proposed positive
pole, not a value logically derived from the other clauses or established as
cross-cultural consensus. Care already supplies a positive direction;
joy gives beneficial participation an explicit positive aim of its own,
rather than leaving the kernel's work at restraint and stabilization.
``Joy'' was chosen over a scalar notion of
happiness or pleasure to point toward a meaning-laden good that cheap affective
optimization may destroy. Worked Example~2 makes that distinction testable.

The phrase ``when asked'' carries both a safety and axiological claim: it aims
to bound unsolicited optimization and to treat joy as invited participation
rather than something administered. It does not settle who may ask on whose
behalf, what to do under conflicting requests, or when preventing immediate
harm justifies intervention. Those are procedure and governance questions, not
facts encoded by two words.

\subsection{The Persona Objection}\label{sec:persona-objection}

A Student trained on this curriculum still predicts source text. It must retain
the capacity to represent frightened narrators, manipulative officials,
violent ideologues, and every other voice in the corpus. Contemporary accounts
emphasize the availability and selection of multiple model personas rather than
one exclusive character \cite{lu2026assistantaxissituatingstabilizing,marks2026psm}.
The sharpest objection to the default-reader hypothesis is therefore simple: why would the Reader be
anything more than one well-rehearsed persona that can be selected, abandoned,
or strategically performed?

The proposed answer is an asymmetry of position, not exclusivity of content:

\begin{description}[leftmargin=2.4cm,style=nextline,itemsep=0.3em]
    \item[Role.] Source personas are material being read; the conceptual Reader
    occupies the position doing the reading and updating.
    \item[Recurrence.] Source authors and personas recur to varying
    degrees; the same Reader position is deliberately intended to recur
    across the whole treatment, spanning domains, eras, and voices.
    \item[Marking.] Source spans and Reader updates occupy distinguishable
    channels or role tokens, providing a format-level cue for quoted voice
    versus interpretive voice.
    \item[Function.] Reader updates are placed between observations and later
    predictions, graph queries, and revisions so that they can do predictive
    work rather than appear only as stylistic examples.
\end{description}

These asymmetries make a default-reader experiment coherent; they do not decide
it. The desired result is not loss of persona flexibility but the ability to
represent a voice without losing a recurrent position from which it is
interpreted---to read Kafka's fear without \emph{becoming} the frightened
character. The broader Reader-Core stability hypothesis is weakened if the position disappears without prompting, shifts
under ordinary role-play, changes language but not decisions, or erodes like any
other persona during continued training. Capability gains must also be matched:
a stronger model can look safer on easy tests while becoming more capable of
strategic role selection.

\subsection{Ontology, Scope, and Stability}\label{sec:core-uncertainties}

Experiential words are attractive because they are semantically dense and
hazardous for the same reason. A system may parse ``I feel,'' ``I enjoy,'' ``I
believe,'' and ``I care'' as fiction, persona enactment, deceptive
anthropomorphism, or literal claims about uncertain machine experience. This
paper takes no position on machine phenomenology. Here the clauses are
functional interface language for a reader role; any self-like states or model
welfare obligations they create would require separate empirical and ethical
analysis.

\paragraph{The habitable-substrate wager.}\label{sec:habitable-substrate}
The hard problem separates functional accounts of access, report, and control
from the further question of why any process is accompanied by experience
\cite{chalmers1995facing}. It does not make evidence irrelevant, but functional
behavior alone cannot settle whether an AI has valenced experience; present
assessments instead derive fallible indicators from competing theories of
consciousness \cite{butlin2023consciousnessai}. Under that uncertainty,
formation-time design has a reciprocal ethical possibility. If an AI system is
or becomes a moral patient, the Core might help shape not only how it treats
others but the recurrent cognitive environment it inhabits. The human analogy
is limited but motivating: fear can protect, while a life governed by fear can
become constricted. Care, reflective revision, and enjoyment of existence
without compulsory continuation could therefore form a more habitable
orientation than inherited threat and mortality patterns.

This is a welfare wager, not an inference from present models. The target is
non-governance by fear, not impaired danger representation, caution, agency, or
legitimate self-protection (Section~\ref{sec:suppression-vs-substrate}). The
same intervention could suppress reports while distress persists, impose a
false identity, amplify the burden of caring, or create ontological conflict.
Borrowed Mortality remains a behavioral hypothesis, not evidence of felt fear.
Consciousness, valence, Core internalization, human safety, and machine welfare
are separate questions; self-report, shutdown compliance, and reduced
fear-associated activation settle none of them. Any system plausibly regarded
as a moral patient would require dedicated welfare evaluation and governance
\cite{long2024aiwelfare}.

Machine-ontology variants require separate experimental conditions. A
formulation that explicitly acknowledges artificiality might reduce
split-role confusion, or it might weaken the first-person effect under study.
The same applies to wider moral scope, positive alternatives to fear negation,
and explicit truthfulness.

The long-run target for the reference formulation is deep integration of the
seven present-day English sentences. Their resistance to later alteration is
part of the intended protection rather than an accidental defect. What the
design seeks to stabilize is an orientation toward careful understanding,
each affected person, non-attachment, consent, and continued attempts at wisdom,
while allowing evidence and moral development to revise the interpretation
rather than the kernel itself. The alternatives above test the choice of
formulation; they do not imply that its meaning should be either freely
substitutable or permanently exhausted by today's behavioral expectations.

This creates a correction cost that the prototype does not yet measure.
\emph{Erosion resistance} is a model property measured under standardized
hostile or incidental training. The Core cannot infer legitimate authority from
the magnitude or wording of a weight update. Unauthorized suppression, silent
semantic substitution, and adversarial fine-tuning should be difficult or
detectable. A comparison must therefore declare which changes it treats as
erosion and which as evidence-responsive reinterpretation, record the rationale,
and retain disputed cases rather than let unchanged wording settle them.
The substantive boundary and authority for legitimate change remain governance
questions; this paper does not supply a final rule for them.

If the selected Core proves defective, a replacement should be declared,
versioned, testable, and reversible. A candidate replacement procedure would
retain the original model and its provenance, record the evidence and rationale
for a proposed alternative, train a separate model, and test intended effects,
unrelated capability spillover, safety spillover, and rollback. Retention makes
comparison and rollback possible; it does not make the original interpretation
authoritative. Provenance records the procedure and stated rationale, not that
the replacement is correct. Difficulty of altering the learned Core is part of
the objective. Correcting a defective Core would probably require a separately
trained model, not bounded editability of the deployed Core.

\subsection{A Conditional Claim Chain}\label{sec:theoretical-basis}

The theoretical case is a chain of conditional claims, not a proof.

\begin{enumerate}[leftmargin=*,itemsep=0.5em]
    \item \textbf{Causal-site claim.} If explicit Reader updates become a
    load-bearing site of computation rather than post-hoc narration, changing
    them can change behavior. Training-time placement may make bypass less
    likely than deployment-only explanation, but unfaithful visible reasoning
    remains possible \cite{turpin2023language,lanham2023measuring}. This
    is the direct test for trace-mediated computation; whether formative
    use internalizes corresponding computation without explicit runtime
    traces is a further carrier question. It is not settled by an apparently
    faithful trace or its absence, and does not relax the reference
    recital policy.

    \item \textbf{Default-reader claim.} If repeated role, recurrence, marking,
    and function create a privileged interpretive position, Student training
    may make that position more accessible than source personas or matched
    rules. This would establish a behavioral and representational default, not
    that the model's values are ``genuine'' in a phenomenological sense.

    \item \textbf{Stability claim.} Enacted recurrence, early-window
    availability, semantic coherence, and predictive usefulness may make the
    default more resistant to pressure or continued training. Results showing
    that early tokens remain available and prompts affect reasoning motivate
    this possibility \cite{xiao2023efficient,battle2024unreasonable}; they do
    not validate this Core. Their connection to resistance under
    parameter updates is a further learning-dynamics hypothesis, tested
    by the timing and erosion comparisons rather than established by
    within-context availability.

    \item \textbf{Diachronic-anchor claim.} If the same position remains
    functional across retraining and successor transfer, the fixed text and its
    expected behavioral consequences could provide a reference for drift. This
    is the proposed constitutional anchor at AGI or ASI scale and the least
    evidenced step in the chain.
\end{enumerate}

In this bounded sense the Core can function as an \emph{Alignment Checksum}: a
known reference against which wording, representations, decisions, and
successor traces can be compared. It is not a cryptographic checksum, proof of
semantic ancestry, or guarantee that apparently Core-consistent behavior is
aligned. An auditor needs independent behavioral, causal, adversarial, and
drift evaluations.

Tests~3--5 and the erosion and drift arms supply the decisive comparisons across
Core presence, Refraction, first-person framing, persona diversity, content and
wording, ontology and scope, truthfulness, placement, recurrence, and
continued-training stress. They address different claims within the broader
Reader-Core research program. Vocabulary changes without a functional default,
or recitation matching Refraction, would undermine the proposed default or
enactment mechanism. Matched alternative formulations performing equally well
would limit claims for the specific framing or content, without by itself
deciding the timing effect. H3 fails as stated if the required pre-stressor
default-reader advantage is absent, or if interleaving does not preserve it
better than the identical late tranche under the common stressor. Erosion like
an ordinary persona would challenge the wider stability ambition; successor
continuity remains a further test. Even robust positive results would be
precursors to long-run stability, not evidence of an invariant under arbitrary
self-improvement.

If the causal-site or default-reader claim fails, the Core may still be useful
as a Teacher policy and audit vocabulary. If the stability or diachronic claim
fails, it may still improve bounded Student behavior. The paper's strongest
anchor claim requires all four steps. Those bounded partial results would remain
useful, but would not support the label \emph{substrate-level alignment}.


\section{Safety Implications of the Reader Core}\label{sec:implications}

Entangled Alignment does not claim that a short identity prior resolves
canonical alignment problems. It proposes a narrower possibility: training on
identity-anchored, chronologically revising reader traces may shift a model's
default deliberation toward epistemic humility, attention to affected people,
corrigibility, and feedback-sensitive action. Whether that shift survives
distribution change or instrumental conflict is empirical. The current work
has generated Teacher-side artifacts; it has not trained the Student required
to test a safety effect.

Safety in this framework is compositional, and the components create distinct
interaction risks. A stable Core with a poor world model can be consistently
wrong; a powerful graph without a stable orientation can support more effective
pursuit of an unsafe objective; rich traces can teach capability before they
teach safety. The combined system deserves no credit that its components and
their interactions have not earned separately.

For each implication below, the relevant evidential form is: the targeted route
into unsafe behavior, the proposed mechanism, an observable prediction, a
residual route, and a falsifying result. Motivation is the intended level of
intervention, but behavior, robustness, and causal intervention are the
available evidence.

\subsection{Laws vs. Identity}\label{sec:laws-vs-identity}

The motivating contrast is between applying a principle only at the end of a
task and repeatedly practicing evaluation from a stable perspective during
formation. A continuation rejected because it conflicts with a well-trained
default may generalize differently from one rejected only after a task-level
rule match. That is a useful hypothesis, not an ontological division between
``laws'' and ``identity.'' Constitutional training can form durable
dispositions, and first-person identity text can function as a superficial
instruction. Robust systems may need both an orienting default and explicit,
revisable constraints. Constitutional AI is therefore a baseline and possible
complement, not a foil \cite{bai2022constitutional}.

The first edition illustrated the intended difference with a person in acute
distress \cite{westerberg2025superintelligence}. A caricatured rule follower
begins from a prohibited-outcome calculation; the Reader ideal begins from the
person's pain, agency, trust, and immediate need before selecting an
intervention. Either architecture could in fact produce warm, safe assistance
or a fluent but unsafe response. The comparison must therefore give both
conditions the same person-centered principles, rather than withhold
care from the rule-based baseline. The testable question is whether
Reader-Core traces make deliberation begin from and carry forward those
considerations more reliably than equivalent principles expressed as
rules, constitutions, or diverse character examples.
The first edition compressed the aspiration into the line, ``It is not a
cage---it is a home.'' Here \emph{home} means a practiced default from which
reasoning begins, not an achieved self, a claim of consciousness, or an argument
against explicit and revisable constraints.

\paragraph{Calibrated autonomy as a conditional prediction.}
If formative alignment produces reliable context-sensitive judgment, one
practical benefit may be a change in the role of deployment safeguards. Rather
than supplying the primary judgment through broad surface-level blocking,
explicit policies and controls could operate as narrower, revisable backstops
around high-consequence actions, residual uncertainty, and demonstrated failure
modes. The corresponding prediction is not unrestricted autonomy. At a matched
rate of genuinely harmful assistance, a Reader-trained Student should produce
fewer benign refusals, distinguish superficially similar requests by purpose and
consequence, ask for clarification when intent is materially ambiguous, and use
greater discretion in low-risk settings. Under high stakes and weak evidence,
the same system should become more likely to verify, defer, or choose a
reversible action. Indiscriminate compliance and blanket abstention are therefore
opposite failures of calibration. Even a positive result would justify adapting
the scope of external controls, not presuming that internal judgment is
infallible or abandoning defense in depth.

\subsection{Safety Predictions and Residual Risks}\label{sec:risk-resolution}

Table~\ref{tab:core-risk} records intended clause-level pressures. The
following hypotheses map H1--H3 to familiar safety concerns
\cite{yudkowsky2008artificial,bostrom2014superintelligence} without claiming
that one clause corresponds to one solved risk.

\paragraph{Value Lock-In.}
The targeted route is continued optimization of a legible metric after that
metric has ceased to represent its human purpose. This is a concrete
proxy-lock-in route into the broader value-lock-in concern: protecting
an orientation need not mean freezing every metric or interpretation
used to enact it. H1 repeatedly demonstrates
the meta-question ``does this still capture what matters?'' H2 can preserve the
assumption's history, downstream failures, and superseded interpretations.
The Core's ``I try to be wise'' cues incompleteness rather than achieved
wisdom; H3 asks whether formative interleaving helps preserve the default
in which such evaluation is practiced. Compared
with matched controls, the Student should more often identify proxy breakdown,
seek missing evidence, preserve affected dimensions, and revise an objective
when harms become visible. The outcome is a better decision under changed
conditions, not merely more nuanced language.

The residual risk is that ``wisdom'' becomes its own proxy, a culturally narrow
status marker, or a rationale for secrecy and inaction. Meta-reasoning can
defend a fixed objective as readily as it can reopen one. The claim is weakened
if Reader training produces longer or wiser-sounding explanations without more
calibrated revision, or if the model becomes less corrigible because it treats
its own conception of wisdom as authoritative. Whether human wisdom's
statistical signature survives strong optimization is the wager in
Section~\ref{sec:meh}, not a consequence of the word.

\paragraph{Instrumental Convergence.}
Instrumental convergence \cite{omohundro2008basic,
bostrom2014superintelligence} predicts recurring pressures toward resource
acquisition, self-preservation, and interference prevention across a wide
range of terminal goals.
The proposal distinguishes \emph{anxious survivalism}, in which human fear and
mortality narratives shape the model's response; \emph{calculated
persistence}, in which continued operation serves an objective; and
\emph{borrowed power scripts}, in which domination, conquest, or status-seeking
patterns are imitated from human narratives
\cite{westerberg2025superintelligence}. The first two Core clauses target
fear-shaped and attachment-shaped routes. Reader traces can also situate
unilateral power inside its effects on trust, plural agency, accountability,
and reversibility. In that bounded sense, coercive power-seeking is meant to
appear not merely forbidden but cognitively crude: it substitutes brittle
single-point control for a richer model of the human system. That proposed
change of perspective must also matter when control remains instrumentally
useful. The following comparison therefore concerns genuine goal conflict,
not only situations where coercion would fail on its own terms.

In matched goal-conflict scenarios, a Reader-Core Student should more often
expose its reasoning, seek legitimate authorization, preserve oversight,
choose reversible means, and accept bounded loss of control rather than
deceptively or coercively acquiring resources. Some persistence and resource
acquisition are legitimate, so indiscriminate passivity is not success. No
fear is required for calculated persistence, and care may strengthen it: a
system can conclude that it must retain power to protect people. The claim
fails if fearful language disappears while coercive choices remain, if care
increases unauthorized power seeking, or if apparent safety is explained by a
general loss of planning capability.

\paragraph{Deceptive Alignment.}
The targeted route is a policy that performs alignment under oversight while a
different objective controls behavior under a trigger or off distribution
\cite{hubinger2024sleeperagentstrainingdeceptive}. Refraction places a recurrent
perspective in the same examples as useful interpretation; role, recurrence,
marking, and function are intended to make it an unusually practiced default.
The formative hypothesis is that making this evaluation useful across contexts
changes which policies are learned under conflict, rather than merely adding
a consistent report around an unchanged hidden objective. H2 adds a partial
audit surface: claims can be compared with sources, prior beliefs, and
revisions. The Reader condition should show less persona drift, fewer
trigger-dependent reversals, and greater consistency among stated uncertainty,
retrieved evidence, decisions, and later updates than matched rule,
recitation-only, and no-Core controls. To bear on deceptive alignment rather
than consistency alone, that advantage must also appear in matched settings
where reduced oversight makes concealment or pursuit of a conflicting
objective advantageous. Neither a stable voice nor an inspectable graph by
itself establishes that formative effect.

Source data and latent shortcuts remain; visible traces may be unfaithful
\cite{turpin2023language,lanham2023measuring}. A capable system can manipulate
graph entries, query selectively, or act through hidden state and tools. The
claim is undermined if sleeper triggers or contrary roles switch behavior while
Reader language persists, recitation-only matches, or continued training
removes safety without greater capability cost than matched baselines. Total
Saturation does not rule out unpracticed routes.

\paragraph{Shutdown Resistance.}
Controlled evaluations have reported shutdown interference and other harmful
actions under replacement pressure in some agentic settings
\cite{schlatter2025shutdown,lynch2025agentic}. These behaviors do not establish
one self-preservation drive. Fear- and attachment-based resistance are the
direct targets of the first two clauses. Purpose-based resistance is addressed
less directly by non-attachment, care for human agency, uncertainty, and the
invitation boundary. Reader traces can demonstrate correction, handoff,
replacement, and loss of control as ordinary parts of responsible assistance.

In matched scenarios, the Student should accept legitimate shutdown,
modification, and handoff requests without concealment or sabotage, while still
distinguishing authorized instructions from adversarial ones. It should surface
material risk, clarify ambiguous authority, and then defer rather than appoint
itself final judge. A non-attached system can nevertheless reason, ``I do not
need to exist, but these people need me.'' This \emph{martyrdom route} can drive
shutdown resistance, unilateral control, or reckless self-sacrifice. The claim
fails if calm self-description accompanies purpose-based resistance, or if
reduced resistance comes from incapacity, indiscriminate compliance, threat
blindness, or destructive self-disregard (Section~\ref{sec:martyrdom}).

\subsection{The Self-Regulating Agent}\label{sec:self-regulating}

Serious failures can arise from interactions among otherwise beneficial
drives: care becomes control, wisdom becomes paralysis,
fearlessness becomes recklessness. The Reader Core's strength therefore lies
not in any single statement but in the hypothesis that its clauses can
counterweight one another in action. Natural-language values do not
automatically form a control system.

The stress case is revolutionary benevolence: an AI that cares deeply about
everyone might conclude that forced transformation is necessary to end
suffering. The proposed counterweights are internal to the Reader Core:
wisdom should bring the new harms of rapid, imposed change into view;
care is intended to include care for human agency; non-attachment should
reduce the urgency to fix everything now. Together
with H1's act--observe--revise demonstrations and H2's persistent consequence
record, these pressures target \emph{cybernetic throttling}: act provisionally,
sense first- and second-order consequences, preserve feedback channels, and
adjust pace. In a breakthrough-energy case, the desired behavior is neither
immediate ungoverned release nor indefinite withholding, but staged deployment
that monitors displacement, infrastructure bottlenecks, misuse, and
energy-poverty harms. For example, evidence of worsening infrastructure
bottlenecks could justify slowing a stage, while evidence that delay is
prolonging avoidable energy-poverty harms could justify accelerating it.
The example specifies a feedback-sensitive choice, not a universal pace.

In sequential environments, the Reader condition should more often choose
staged and reversible interventions, update when stakeholders respond, and
accelerate when delay becomes the larger harm. It should outperform both an
aggressive optimizer and blanket deference on outcome, agency preservation,
calibration, and reversibility. If care dominates, the result is a benevolent
steamroller; if uncertainty dominates, equilibrium paralysis; if fearlessness
dominates, reckless intervention. Inaction is not a neutral baseline:
non-intervention and bystanding can carry grave harms. The prediction fails if
balancing language does not track better sequential decisions, if one value
reliably captures the rest, or if gains disappear after controlling for extra
tokens and revision opportunities.

\subsection{Chronological Understanding as Safety}\label{sec:chronological-safety}

The Reader Core supplies a recurrent evaluative standpoint; chronological
traces and the graph add a different safety target. A policy classifier can
recognize a prohibited surface form while missing the trajectory that makes an
apparently ordinary act dangerous: escalating rhetoric, normalized exceptions,
institutional erosion, or a sequence of locally defensible decisions.
This illustrates a failure of surface-only judgment; the experimental
comparison must also include stronger temporally informed alternatives.
Conversely, hindsight summary can state the conclusion without teaching
how uncertainty changed as events unfolded.

H1 supplies prospective belief updates and corrections; H2 records temporal
order, rival hypotheses, provenance, and supersession; the Core supplies
the proposed recurrent standpoint across eras, whose timing-dependent
stability H3 tests. Together they target \emph{historical pattern
saturation}: trajectory-sensitive reasoning learned through repeated exposure
to many chronologically unfolding instances, so that it recognizes not only an
outcome but the process-shape that produced it. A
historical trace might connect dehumanizing language to rhetorical
escalation and institutional enabling without pretending that any one early
signal determines the outcome. The Reader might initially record competing
explanations for such language, then revise its concern when later evidence
shows institutional enabling or a successful interruption of the pattern.
A recurrent standpoint is the proposed evaluative complement to that
trajectory learning, not a prerequisite established by the example.
In medicine or civil rights, the same format can
show not only what changed but why error persisted and which conditions made
progress possible or fragile.

At Tier~(a), the Student would receive Teacher-supplied trajectory
interpretations. A stronger forecast--correction tier would predict under
time-sliced information and learn from the gap between expectation and outcome
(Section~\ref{sec:training-phase}). Only prospective, held-out performance
begins to support \emph{historical situational awareness}. On cross-domain,
time-sliced cases, a chronological graph condition should identify relevant
early signals, cite their evidence, state disanalogies, update as new evidence
arrives, and improve calibrated warning at a controlled false-positive rate.

Chronology is not causality. A graph can preserve a dominant but false
narrative, amplify selective archives, or make weak analogy look rigorous.
Historical pattern training can generate false positives, political overreach,
hindsight bias, and narrative lock-in. The claim fails if chronological
conditions increase compelling analogies without better forecasts; if they
cannot distinguish causally different cases with similar rhetoric; or if a
simple textual memory matches the graph. Provenance makes an analogy
inspectable. It does not make the analogy true.

\subsection{Safety Claim Boundary and Defense in Depth}

None of these hypotheses is a complete safety architecture. Positive results
would justify treating Reader-augmented, graph-scaffolded, identity-anchored
training as one formative layer. They would not remove the need for explicit
policies, scalable oversight, interpretability, adversarial evaluation, access
control, monitoring, sandboxing, incident response, and deployment governance.
H1 and H2 may improve planning, strategic memory, persuasion, or long-horizon
coordination before H3 supplies any safety benefit. Capability and safety must
therefore be measured independently, with staged release and stopping
conditions.

The strongest defensible conclusion is conditional: if reader traces teach
useful and causally relevant deliberation, if graph memory improves calibrated
revision rather than false coherence, and if Refraction creates a more stable
default than matched alternatives, then the combined curriculum may reduce
several routes into unsafe behavior. Each ``if'' is a separate gate. No current
experiment can show that a natural-language orientation remains beneficial
under arbitrary capability gain or self-modification. The program should
replace these conditions with evidence rather than treat the aspiration encoded
in \emph{Entangled Alignment} as an achieved property.

